\documentclass[../main.tex]{subfiles}
\begin{document}

\FloatBarrier
\section{Overview}
\label{sec:hybrid-overview}

This chapter reports the experimental results, centred on the proposed hybrid model of Section~\ref{sec:hybrid}. Two questions are put to it: whether its best configuration improves on the single-encoder reference in-domain, and whether the architecture and dataset combination reinforce each other out of domain, where deployment actually happens.

Two findings from the single-encoder baseline study frame the comparison and are established in Appendix~\ref{chap:baseline}. First, the strongest baseline reaches QWK 0.537, the reference level the proposed model must exceed. Second, the OpenFace and I3D feature families, though they score almost the same, fail on different clips, so a model that learns to combine complementary streams should improve on any single encoder. That study also fixes the two protocol choices this chapter inherits, cross-entropy as the development loss and QWK as the single selection metric, and justifies them empirically.

The chapter proceeds as a sequence of experiments, each stating its aim, its setting, its result, and the conclusion that leads to the next. Section~\ref{sec:setup-hybrid} states how the experiments are conducted, including the loss forms, the training regime, and the 3$\times$3 evaluation matrix that every result in the thesis uses.

\FloatBarrier
\section{Experimental Setup}
\label{sec:setup-hybrid}

The experiments evaluate one architecture family, the proposed hybrid model of Section~\ref{sec:hybrid}, over its full design space. All $3^{5}=243$ per-group encoder assignments are instantiated, each with the I3D stream enabled and disabled, giving 486 configurations. Every configuration is trained with cross-entropy, the development loss established in the baseline study (Appendix~\ref{chap:baseline}, Section~\ref{sec:loss-tradeoff}), on each of the three sources and evaluated under the training regime and 3$\times$3 protocol defined below. Fixing the loss isolates the effect of the architecture.

The reporting follows one metric policy throughout: QWK selects every best model, macro-accuracy and macro-MAE accompany the headline results, and accuracy appears only as the cautionary metric, for the last time in Section~\ref{sec:hybrid-crossdomain}. The empirical case for this policy is made on the baselines in Appendix~\ref{chap:baseline}. A headline comparison reports the tuned maximum of QWK over the design space; a claim about the family is reported as a mean and a distribution over the 243 or 486 configurations, so a family-level effect is robustness across the design space rather than a single tuned point.

The experiments follow the four research questions. Section~\ref{sec:ablation} confirms that the evaluation protocol carries over to this larger population (RQ1's frame condition); Section~\ref{sec:group-marginal} isolates which group's encoder choice matters (RQ2); Section~\ref{sec:i3d-effect} measures I3D fusion with a paired ablation (RQ3); Section~\ref{sec:error} pins down the in-domain comparison with the baseline (RQ1); Section~\ref{sec:private} takes the model out of domain and to the private set (RQ4); and Section~\ref{sec:gen-discussion} discusses why the design works.

\subsection{Loss Functions}
\label{sec:losses-method}

Three loss families are used (Section~\ref{sec:losses-lit}): cross-entropy, weighted cross-entropy, and the ordinal squared-EMD loss of Equation~\eqref{eq:emd}. The baseline study sweeps all three to measure the effect of the loss (Appendix~\ref{chap:baseline}); the proposed model is trained only with the development loss selected there. The class weights for the weighted and ordinal variants are inverse-frequency weights computed on the training split and normalised to mean~1: for class $c$ with $n_c$ training samples and $C$ non-empty classes,
\begin{equation}
w_c = \frac{1}{n_c}\cdot\frac{C}{\sum_{k} \frac{1}{n_k}},
\qquad\text{so that}\quad \frac{1}{C}\sum_c w_c = 1 .
\end{equation}
The ordinal loss is applied in its class-weighted form as well, multiplying each sample's CDF distance by $w_y$.

For the proposed hybrid model the objective also includes the auxiliary term. With main logits $z$ and per-stream auxiliary logits $\{a^{(s)}\}$, the total loss is
\begin{equation}
\mathcal{L} = \mathcal{L}_{\text{main}}(z,y) \;+\; \lambda_{\text{aux}}\cdot\frac{1}{S}\sum_{s=1}^{S}
              \mathcal{L}_{\text{CE}}\!\bigl(a^{(s)},y\bigr),
\end{equation}
where $\mathcal{L}_{\text{main}}$ is the main-head loss, $\lambda_{\text{aux}}=0.2$ is the auxiliary weight, and the auxiliary heads are always trained with plain cross-entropy over the $S$ streams (five groups, or six with the I3D stream enabled). The formulation admits any of the three losses for the main head; the ablation here fixes $\mathcal{L}_{\text{main}}$ to cross-entropy, so its results isolate the effect of the architecture from that of the loss.

\subsection{Training Protocol}
\label{sec:training}

All models are optimised with Adam~\cite{adam} at a learning rate of $1\times10^{-4}$. Training uses mini-batches (128 for the baselines, 64 for the ablation) and early stopping on the validation loss (the CMOSE \texttt{unlabel} split and its analogues) with patience 10 and an epoch cap (800 for baselines, 400 for the ablation). The checkpoint with the lowest validation loss is restored for testing. Following the convention adopted in this project, only the training and validation loss are recorded per epoch; the final metrics are evaluated once, on the held-out test split. Runs use a fixed seed (42) and optional mixed precision on GPU, and every run converges stably and early-stops well before the cap, so the reported differences are architectural rather than optimisation artefacts.

\subsection{Metrics and the 3$\times$3 Evaluation Matrix}
\label{sec:eval}

Generalisation is measured with a 3$\times$3 matrix: every model and loss is trained on each of \{CMOSE, DAiSEE, Combined\} and evaluated on each of \{CMOSE-test, DAiSEE-test, Private\}. The diagonal cells are the in-domain results; the off-diagonal cells and the Private column measure transfer to unseen distributions.

Each evaluation reports the six metrics of Section~\ref{sec:metrics-lit} (accuracy, macro-accuracy, MAE, macro-MAE, Cohen's kappa, and QWK), but QWK is the single primary metric: it alone selects the best model. The other five are secondary, reported for analysis and comparability but never used to pick a winner. The empirical justification for selecting on QWK alone is given in the baseline study (Appendix~\ref{chap:baseline}, Section~\ref{sec:protocol}). Row-normalised confusion matrices and per-class F1 scores are used for error analysis.

The study comprises roughly 1{,}500 trained runs (the five baselines $\times$ three losses $\times$ three sources, and the 243 per-group configurations of the proposed hybrid model with the I3D stream enabled and disabled); all are aggregated programmatically, and every figure and table of this chapter and Appendix~\ref{chap:baseline} is produced from the aggregate by a single command.

\FloatBarrier
\section{Evaluation Properties of the Proposed Model}
\label{sec:ablation}

The evaluation protocol of Section~\ref{sec:eval} was fixed on fifteen single-encoder baselines (Appendix~\ref{chap:baseline}); this experiment confirms that its three properties hold for 486 configurations of a different architecture family, none of which existed when the protocol was set. Figure~\ref{fig:hybrid_ablation_all_metrics} plots the distribution of the two variants of the proposed hybrid model (I3D stream enabled and disabled) on every metric, each panel against the best baseline on that metric, so all six metrics appear here for reference. Three properties reproduce.

\begin{figure}[ht]
\centering
\includegraphics[width=\linewidth]{hybrid_ablation_all_metrics.png}
\caption{Ablation of the proposed hybrid model across all 243 per-group encoder
configurations, with the I3D stream enabled and disabled, on all six metrics (in-domain
CMOSE). The dashed line is the QWK-selected baseline (OpenFace TCN, cross-entropy).}
\label{fig:hybrid_ablation_all_metrics}
\end{figure}

The metric structure reproduces. The Kendall rank correlation between the six metrics over the 486 in-domain configurations gives the same two-block picture as the baseline study (Appendix~\ref{chap:baseline}, Section~\ref{sec:protocol-metrics}), a micro block, a macro block, weak cross-block agreement, and QWK again the only bridge, which reflects the metrics on this task, not any model family. Accuracy remains insensitive: every configuration falls within a narrow accuracy range around the baseline's 0.761 (top-left panel), so a model chosen on accuracy would see no architecture signal, while the QWK panel separates the two variants cleanly. The cross-domain collapse reproduces too: taken out of domain (Section~\ref{sec:hybrid-crossdomain}), the proposed model's matrix has the same shape as the baseline matrix, a healthy diagonal (0.605 on CMOSE) and an off-diagonal collapse, with accuracy again masking it.

All three properties thus hold across a 32-fold larger model population, so the protocol carries over without amendment (RQ1's frame condition), and the experiments that follow can use it to judge the architecture.

\FloatBarrier
\section{Encoder Choice per Group}
\label{sec:group-marginal}

The design space varies five group$\rightarrow$encoder choices at once; the aim here is to find which group's choice actually moves QWK (RQ2), and whether any preference is specific to one dataset or preserved under shift. Table~\ref{tab:T5_1_group_marginal_combined} marginalises each group's encoder choice in both regimes, in-domain and over the unseen-target cells.
\begin{table}[ht]
\centering
\caption{Per-group marginal effect of encoder choice on mean QWK, in-domain (CMOSE) and pooled over the unseen-target cells.}
\label{tab:T5_1_group_marginal_combined}
\fittable{%
\begin{tabular}{llrrrrl}
\toprule
Group & Regime & TCN & Transformer & LSTM & Spread & Prefers \\
\midrule
Gaze & In-domain & 0.539 & 0.537 & 0.534 & 0.006 & TCN \\
Gaze & Unseen targets & 0.093 & 0.095 & 0.094 & 0.002 & Transformer \\
Eye landmarks & In-domain & 0.538 & 0.537 & 0.535 & 0.003 & TCN \\
Eye landmarks & Unseen targets & 0.096 & 0.092 & 0.094 & 0.004 & TCN \\
Face landmarks & In-domain & 0.537 & 0.534 & 0.539 & 0.006 & LSTM \\
Face landmarks & Unseen targets & 0.092 & 0.096 & 0.093 & 0.004 & Transformer \\
Head pose & In-domain & 0.545 & 0.523 & 0.541 & 0.022 & TCN \\
Head pose & Unseen targets & 0.100 & 0.091 & 0.090 & 0.010 & TCN \\
Action units & In-domain & 0.536 & 0.537 & 0.537 & 0.001 & Transformer \\
Action units & Unseen targets & 0.093 & 0.096 & 0.093 & 0.003 & Transformer \\
\bottomrule
\end{tabular}%
}
\end{table}


Only head pose shows a clear preference. In-domain it prefers a TCN (mean QWK 0.545, ahead of LSTM 0.541 and Transformer 0.523, a spread of 0.022), and macro-MAE agrees, lowest under the TCN; the other four groups are insensitive to the encoder (QWK spread $\le0.006$). Under shift head pose again has the largest spread (0.010) and again prefers the TCN, while the other groups' spreads (0.002--0.004) flip winners arbitrarily, which is noise. The consistency across regimes shows the head-pose rule reflects the signal rather than a tuning accident: head dynamics such as nodding and turning away are temporally local and suit convolution (RQ2). The design space is otherwise flat, so most of the architecture's benefit must come from decomposing and fusing rather than from the encoder search, which the next experiment measures.

\FloatBarrier
\section{Effect of I3D Fusion}
\label{sec:i3d-effect}

The baseline study predicted, from its error-overlap analysis (Appendix~\ref{chap:baseline}, Section~\ref{sec:openface-vs-i3d}), that fusing the I3D stream should help; this experiment tests the prediction. The design is paired: each I3D-stream-disabled configuration is matched with its exact I3D-stream-enabled twin (same five encoders), so the distribution of paired QWK differences is the I3D effect free of the configuration as a confounder, measured in three regimes (Table~\ref{tab:T5_2_i3d_fusion_effect}).
\begin{table}[ht]
\centering
\caption{Paired effect of the I3D stream on QWK, by evaluation regime.}
\label{tab:T5_2_i3d_fusion_effect}
\fittable{%
\begin{tabular}{lrlr}
\toprule
Regime & Pairs (n) & Mean $\Delta$QWK & Pairs improving (\%) \\
\midrule
Seen target & 972 & +0.060 & 94 \\
Cross-corpus (public) & 486 & -0.021 & 26 \\
Private (unseen) & 729 & +0.005 & 52 \\
\bottomrule
\end{tabular}%
}
\end{table}


Where the test dataset is seen in training, fusion helps. Over the seen-target cells the mean paired gain is $+0.060$ QWK, and 94\% of the 972 pairs improve. In-domain on CMOSE the I3D-stream-enabled configurations are centred above the I3D-stream-disabled ones (median QWK 0.553 against 0.522, visible in Figure~\ref{fig:hybrid_ablation_all_metrics}), and 82\% of the 243 enabled configurations exceed the 0.537 baseline reference, against 26\% of the disabled ones, so the gain belongs to the whole family (RQ3). Beyond the training datasets the gain disappears: cross-dataset the mean paired effect is $-0.021$ QWK with only 26\% of pairs improving, and on the private set it is $+0.005$ with 52\% improving and high variance. I3D's global appearance features carry dataset bias along, the same mechanism behind the I3D MLP's fall from second to fourth under shift (Appendix~\ref{chap:baseline}, Section~\ref{sec:crossdomain-matrix}). The rule (RQ3) is to enable the I3D stream when the target domain is represented in training and rely on the OpenFace groups otherwise; it is a rule about the family average, since the best single private-set model does have the I3D stream enabled but sits in one tail of a wide, zero-centred distribution.

\FloatBarrier
\section{Comparison with the Baseline}
\label{sec:error}

This experiment settles the in-domain comparison (RQ1): does the best configuration of the proposed model beat the best baseline (Appendix~\ref{chap:baseline}), and where does any gain come from. Table~\ref{tab:T5_3_hybrid_topk} lists the top-5 configurations by QWK.
\begin{table}[ht]
\centering
\caption{Top-5 hybrid configs in-domain (CMOSE), sorted by QWK (macro-MAE lower is better).}
\label{tab:T5_3_hybrid_topk}
\fittable{%
\begin{tabular}{llrrrr}
\toprule
Variant & Arch (gaze\_eye\_face\_head\_au) & QWK & Macro-Acc & Macro-MAE & Accuracy \\
\midrule
Hybrid + I3D & TCN\_T\_TCN\_LSTM\_T & 0.605 & 0.582 & 0.489 & 0.784 \\
Hybrid + I3D & T\_TCN\_TCN\_LSTM\_T & 0.591 & 0.606 & 0.483 & 0.787 \\
Hybrid + I3D & T\_T\_LSTM\_LSTM\_TCN & 0.591 & 0.563 & 0.505 & 0.781 \\
Hybrid + I3D & T\_TCN\_T\_TCN\_T & 0.588 & 0.566 & 0.531 & 0.789 \\
Hybrid + I3D & LSTM\_LSTM\_TCN\_LSTM\_T & 0.588 & 0.567 & 0.517 & 0.787 \\
\bottomrule
\end{tabular}%
}
\end{table}


The best configuration improves on the baseline. It has the I3D stream enabled, \texttt{TCN\_T\_TCN\_LSTM\_T}, and every top configuration mixes at least two encoder types. On QWK, the selection metric, the best proposed-model configuration improves on the best baseline by $+0.068$ (0.605 versus 0.537), and the balanced secondaries corroborate it (macro-accuracy 0.582 versus 0.535, macro-MAE 0.489 versus 0.547), so the gain is real but modest.

The per-class recalls localise the gain. Relative to the best baseline, the proposed model improves recall on the ordinal extreme E3 (Highly\nobreakdash-Engage) ($0.38\to0.51$) and on E1 (Disengage) ($0.45\to0.55$) while keeping E2 (Engage) near 0.90, and trades away a little of the rarest class E0 (Highly\nobreakdash-Disengage) ($0.40\to0.37$), reallocating capacity toward the more frequent extreme. The in-domain improvement is therefore real but uneven, clearing the baseline on QWK and the balanced secondaries by strengthening the well-populated upper-engagement region and the middle disengaged class, while the $\sim$2--3\%-frequency E0 class stays the main open challenge (RQ1). The next experiment shows the picture becomes more severe out of domain, and that the architecture's advantage over the baseline widens there.

\FloatBarrier
\section{Cross-Dataset Generalisation}
\label{sec:private}

The aim of the final experiment is the deployment question: how the proposed model transfers to unseen distributions, and whether the architecture gain and dataset combination reinforce each other there (RQ4). It proceeds from the 3$\times$3 matrix to the private set, then asks whether the in-domain leaderboard predicts transfer.

\subsection{Cross-Dataset Matrix}
\label{sec:hybrid-crossdomain}

The proposed model's cross-dataset heatmap (Figure~\ref{fig:crossdomain_hybrid}) shows the same pattern as the baseline matrix (Appendix~\ref{chap:baseline}, Section~\ref{sec:crossdomain}): a healthy diagonal and an off-diagonal collapse of QWK and the balanced secondaries alike, so the proposed model does not rescue cross-dataset transfer. Accuracy stays misleading, for the last time in the thesis's analysis: on DAiSEE$\rightarrow$CMOSE the best QWK any of the 486 runs reaches is 0.10, yet those configurations reach accuracies up to 0.69. A second architecture family thus lands in the same place, so the metric frame is not an artefact of the single-encoder models.

\begin{figure}[ht]
\centering
\includegraphics[width=\linewidth]{crossdomain_hybrid.png}
\caption{Cross-domain generalisation, best proposed-model configuration per cell (selected
by QWK), on QWK, macro-accuracy, macro-MAE, and accuracy.}
\label{fig:crossdomain_hybrid}
\end{figure}

Within the collapse, the proposed model is strongest in every cell. Figure~\ref{fig:crossdomain_delta} subtracts the baseline matrix (Figure~\ref{fig:crossdomain_base}) from the proposed model's matrix cell by cell: the best proposed-model configuration outperforms the best baseline in all nine train$\times$test cells, every cell of the $\Delta$QWK matrix positive. The gains are smallest on the in-domain diagonal ($+0.068$ on CMOSE) and largest on the Private column ($+0.118$, $+0.050$, $+0.094$ from the three training sources). The improvement is a uniform shift of the whole matrix, not a single tuned point, the strongest form the family-level claim can take. Read column by column, the figure ties the two headline comparisons together: its CMOSE diagonal cell is the in-domain gain of Section~\ref{sec:error}, and its Private column is the private-set result developed next.

\begin{figure}[ht]
\centering
\includegraphics[width=0.66\linewidth]{crossdomain_delta.png}
\caption{Cell-by-cell advantage of the proposed model: best proposed-model minus best
baseline QWK over the 3$\times$3 train$\times$test matrix (Figure~\ref{fig:crossdomain_hybrid} minus
Figure~\ref{fig:crossdomain_base}).}
\label{fig:crossdomain_delta}
\end{figure}

\subsection{The Private Set}
\label{sec:private-results}

The main generalisation test is the private set: the 366 clips collected and hand-labelled for this thesis (Section~\ref{sec:private-collection}), processed through the identical OpenFace/I3D pipeline and used for testing only. Because no model is trained on it, it is a test-only, in-the-wild set, with its own recording conditions and label prior (58\% E2, 31\% E3, 2.7\% E0), and the only design lever is the training source. The Private column of Figure~\ref{fig:crossdomain_delta} and Table~\ref{tab:T5_4_private_by_source} report the best baseline and best proposed-model configuration per training source, where two effects combine.
\begin{table}[ht]
\centering
\caption{Private set (test-only): best baseline vs best proposed hybrid model by training source.}
\label{tab:T5_4_private_by_source}
\fittable{%
\begin{tabular}{llrrlrrr}
\toprule
Train source & Best baseline (model/loss) & Baseline QWK & Baseline macro-MAE & Best proposed model (arch) & Proposed QWK & Proposed macro-acc & Proposed macro-MAE \\
\midrule
CMOSE & openface\_transformer/ce & 0.190 & 0.750 & I3D stream disabled, TCN\_T\_LSTM\_TCN\_T & 0.309 & 0.385 & 0.914 \\
DAiSEE & i3d\_mlp/ce & 0.256 & 1.031 & I3D stream enabled, LSTM\_LSTM\_TCN\_LSTM\_LSTM & 0.306 & 0.350 & 0.966 \\
Combined & openface\_transformer/ce & 0.285 & 0.997 & I3D stream enabled, T\_T\_T\_LSTM\_T & 0.379 & 0.385 & 0.877 \\
\bottomrule
\end{tabular}%
}
\end{table}


\begin{enumerate}
  \item \textbf{Training source matters, and the combined dataset performs best.}
  For both model families, training on Combined generalises best to the private set (best baseline QWK: Combined 0.285 versus DAiSEE 0.256 versus CMOSE 0.190): combining mixed source datasets is the most effective simple route to an unseen target, as the baseline study showed (Appendix~\ref{chap:baseline}).
  \item \textbf{The proposed model generalises better, and helps more here than in-domain.}
  It outperforms the best baseline on QWK for every training source, and the gap is largest where it matters: under combined training, QWK 0.379 versus 0.285, a $+0.094$ improvement, wider than the $+0.068$ in-domain gain, so distribution shift amplifies rather than erodes the architecture's advantage. The best private-set model overall (selected on QWK alone) is the combined-trained proposed hybrid model with the I3D stream enabled, \texttt{T\_T\_T\_LSTM\_T}, at QWK 0.379, and the secondary macro-MAE agrees, the lowest of any model (0.877 against the best baseline's 0.997). One caveat in Table~\ref{tab:T5_4_private_by_source}: the CMOSE-only baseline shows a deceptively low macro-MAE (0.75) despite the weakest QWK, by predicting conservatively around the majority classes, exactly the misleading reassurance that selecting on QWK alone is designed to avoid.
\end{enumerate}

The confusion matrices localise where the gain comes from. Figure~\ref{fig:private_confusion_combined} compares the two combined-trained models on the private set: the proposed model captures the ordinal structure better, E2 recall rising from 0.46 to 0.72 with errors concentrated on the diagonal and its neighbours, where the baseline scatters E2 mass into E3 (0.44). Read as per-class F1, the same signature as in-domain (Section~\ref{sec:error}) recurs in sharper form: the proposed model lifts F1 on the well-populated E2 ($0.52\to0.70$) and E3 ($0.49\to0.56$) classes and, where the baseline is all but flat, gives the middle E1 class real mass ($0.09\to0.28$), while the rarest E0 class is traded all the way to zero ($0.07\to0.00$; all ten private E0 clips are misread). Out of domain the model concentrates its limited capacity on the engaged region even more sharply, and the rare disengaged extreme remains the thesis's open failure mode (RQ4).

\begin{figure}[ht]
\centering
\includegraphics[width=\linewidth]{private_confusion_combined.png}
\caption{Private-set confusion of the combined-trained best baseline versus best
proposed-model configuration (row-normalised).}
\label{fig:private_confusion_combined}
\end{figure}

\subsection{In-Domain versus Transfer}
\label{sec:hybrid-indomain-transfer}

The last question repeats, on the proposed model, the baseline test of whether in-domain ability predicts transfer (Appendix~\ref{chap:baseline}, Section~\ref{sec:indomain-vs-transfer}). Each configuration's in-domain cell (trained and tested on CMOSE) is paired with its deployment cell (trained on Combined, tested on the private set), one point per configuration, and Figure~\ref{fig:indomain_vs_generalization_hybrid} overlays the 486 proposed-model configurations on the baseline scatter.

\begin{figure}[ht]
\centering
\includegraphics[width=\linewidth]{indomain_vs_generalization_hybrid.png}
\caption{In-domain QWK (trained and tested on CMOSE) versus transfer QWK (trained on
Combined, tested on the held-out private set); each configuration is one point pairing its
own two runs. Spearman $\rho=0.27$ for the 486 proposed-model configurations and $0.31$
for the 15 baselines.}
\label{fig:indomain_vs_generalization_hybrid}
\end{figure}

The correlation stays weak ($\rho=0.27$, significant only because $n=486$), in the same range as the baseline relationship ($\rho=0.31$): a configuration's in-domain CMOSE ability is, at most, a weak predictor of how it transfers. The architecture's contribution shows as position rather than slope: the proposed model's cloud lies above and to the right of the baseline configurations, reaching both higher in-domain QWK and higher private-set transfer, so the feature-group design lifts both ends together, while combined training remains the lever that places a model in that upper-right region. The private-set winner is therefore found by family-level robustness and combined training, not by reading the in-domain leaderboard.

\FloatBarrier
\section{Discussion}
\label{sec:gen-discussion}

The design choices work for reasons visible in the results. Combining datasets improves generalisation by exposing the model to a more diverse set of distributions, reducing its tendency to overfit to either dataset. The proposed model's margin widens off-domain for the reason it was designed: each behavioural channel is encoded separately and forced to be discriminative through its auxiliary head, so the representation depends less on the distinctive features of any one dataset. The two feature families diverge under shift because the OpenFace descriptors live in subject-centred physical units that mean the same thing in any setup, whereas the I3D embedding encodes global appearance, what changes between datasets, so fusion adds signal only where the test statistics are represented in training.

The main result is a clear but moderate improvement. QWK 0.379 exceeds chance, every single-source model, and the best baseline ($+0.094$), but in absolute terms it is moderate: with macro-MAE 0.877 the model captures the broad engaged-versus-disengaged trend yet is unreliable at fine distinctions, and the rarest class E0 fails in both regimes. The signal is usable as an aggregate, classroom-level measure but not to flag an individual disengaged learner without human confirmation. The open problem is one of data and objective, since no encoder assignment can recover a class the training distribution barely contains, which is why the remedies for imbalance in Chapter~\ref{chap:conclusions} are the natural next step.

\FloatBarrier
\section{Chapter Summary}
\label{sec:hybrid-summary}

The evaluation protocol held under the 486-configuration robustness check: the block structure of the metrics, the insensitivity of accuracy, and the cross-domain collapse all reproduced as task properties. Only head pose has a clear encoder preference (a TCN), so the benefit comes from decomposing and fusing rather than the encoder search (RQ2). I3D fusion is positive where the test dataset is seen in training, neutral-to-harmful beyond it (RQ3). In-domain the best proposed-model configuration outperforms the best baseline on QWK, the single selection metric (0.605 versus 0.537), with the secondary metrics agreeing (RQ1); out of domain it is strongest in every cell of the matrix, and reaches the best private-set model when trained on the combined dataset (QWK 0.379, $+0.094$ over the baseline, wider than in-domain), transferring better than the baselines at every in-domain level (RQ4), with E0 the main failure mode in both regimes. Chapter~\ref{chap:conclusions} draws the findings together, revisits the research questions, and outlines future work.

\end{document}
